\section{Conditioning a relative heat-flux accuracy claim}
\label{sec:conditioning}
The source's Eqs.~(3)--(4) and (B1)--(B2) are correct. For $n_1,n_2>0$, put
$n=n_1+n_2$, $w_s=n_s/n$, $\Delta u=u_1-u_2$ and
$T_s=p_s/n_s$. Recentring each population on the aggregate mean gives
\begin{align}
 p_{ij}&=p_{1ij}+p_{2ij}+mnw_1w_2\Delta u_i\Delta u_j,\nonumber\\
 q_{ijk}&=q_{1ijk}+q_{2ijk}
 +nw_1w_2\bigl(\Delta u_i\Delta T_{jk}
             +\Delta u_j\Delta T_{ik}+\Delta u_k\Delta T_{ij}\bigr)\nonumber\\
 &\hspace{15mm}
 -mnw_1w_2(w_1-w_2)\Delta u_i\Delta u_j\Delta u_k.
 \label{eq:mixture-correct}
\end{align}
The source explicitly defines the unnormalized three-term symmetrization,
so its tensor formula has the required scalar factor three. Direct positive
finite-mixture calculations, including nonunit mass, verify all entries.

For the scalar $xxx$ component and $\Delta u\ne0$, define
$\theta_s=T_{sxx}/[m(\Delta u)^2]$. Its counter-streaming contribution is
\begin{equation}
 q_{\rm CS}=mn(\Delta u)^3C,\qquad
 C=w_1w_2\{3(\theta_1-\theta_2)-(w_1-w_2)\}.
 \label{eq:dimensionless-cs}
\end{equation}
Thus the source's cubic-drift explanation of the empirical parameter collapse
requires qualification. Equation~\eqref{eq:dimensionless-cs} also contains a
temperature-contrast term linear in drift and can vanish by symmetry or
cancellation.

\begin{proposition}[A sufficient conditioned relative-error bound]
\label{prop:conditioned-relative-bound}
Suppose $|q_{sxxx}|\le mn_s a_s v_{{\rm th},s}^3$, with
$v_{{\rm th},s}=(T_{sxx}/m)^{1/2}$ and stated nonnegative bounds $a_s$. Set
$E=\sum_s w_sa_s(v_{{\rm th},s}/|\Delta u|)^3$.
If $E<|C|$, then
\begin{equation}
 \frac{|q_{xxx}-q_{\rm CS}|}{|q_{xxx}|}
 \le\frac{E}{|C|-E}.
 \label{eq:conditioned-bound}
\end{equation}
\end{proposition}
\begin{proof}
The exact residual is $r=q_{1xxx}+q_{2xxx}$, so
$|r|\le mn|\Delta u|^3E$. The reverse triangle inequality gives
$|q_{xxx}|\ge mn|\Delta u|^3(|C|-E)>0$.
Division gives \eqref{eq:conditioned-bound}.
\end{proof}

The bound is sufficient rather than necessary. Its imbalance/cancellation
factor cannot be dropped. Equal population densities and temperatures give
$q_{\rm CS}=0$ at any drift. With equal densities but unequal temperatures,
$q_{\rm CS}=3n\Delta u\,\Delta T_{xx}/4$, which is linear in drift.

\paragraph{A near-Maxwellian counterexample to a speed-ratio-only guarantee.}
Let $W$ take values $-1/\sqrt2$ and $\sqrt2$ with probabilities $2/3$ and $1/3$.
Then $\E{W}=0$, $\E{W^2}=1$ and $\E{W^3}=1/\sqrt2$.
For $\eta>0$, let $S=(W+\eta G)/\sqrt{1+\eta^2}$ with $G$ a standard Gaussian.
Let $H=(S,G_y,G_z)$, with independent standard Gaussian transverse components,
and define $X_\epsilon$ to have law
$(1-\epsilon)\mathcal N(0,I)+\epsilon\mathcal L(H)$, $0<\epsilon<1$.
It has zero mean, covariance $I$, and third moment
$\E{X_{\epsilon,x}^3}=\epsilon/[\sqrt2(1+\eta^2)^{3/2}]$.
Its total-variation distance from $\mathcal N(0,I)$ is at most $\epsilon$.
Give the populations equal densities and velocity laws
\begin{equation}
 V_1=Ue_x+\sigma X_\epsilon,\qquad
 V_2=-Ue_x+\sigma X_\epsilon.
 \label{eq:smooth-skew-mixtures}
\end{equation}
Each population density is a positive smooth Gaussian mixture, within
$\epsilon$ in total variation of its matching Maxwellian. Their intrinsic
temperature tensors are both $m\sigma^2I$, so
$q_{\rm CS}=0$, whereas
$q_{xxx}=mn\sigma^3\epsilon/[\sqrt2(1+\eta^2)^{3/2}]\ne0$.
Consequently the relative closure error is one even when both $\epsilon$ and
the thermal-to-beam speed ratio $\sigma/U$ are arbitrarily small.
The absolute residual does tend to zero relative to $mnU^3$.
The recorded calculations use $\epsilon=0.01$ and integrate the smooth
laws with positive Gauss--Hermite quadrature.

This construction does not disprove the source's measured collapse in its
particular family of reconnection simulations. It identifies assumptions
needed to extend that empirical observation to a relative-error guarantee.
The source scan uses initial speeds; inserting those into a local bound would
also require estimates of their subsequent evolution.

\paragraph{The valid limitation of a memoryless one-fluid closure.}
Additional population information can be genuinely necessary. For example,
let $n=m=1$, and let each population be an isotropic Gaussian of variance
$3I/4$. Choose a weight $w$, set
$\Delta=[4w(1-w)]^{-1/2}$, and assign means
$(1-w)\Delta e_x$ and $-w\Delta e_x$.
For both $w=1/2$ and $w=1/4$, the aggregate density is one, the mean is zero,
and the pressure is $\diag(1,3/4,3/4)$. The respective $q_{xxx}$ values are
zero and $1/(4\sqrt3)$. Hence no arbitrary-distribution algebraic function
of only $n,u,p$ can recover both third moments. This supports the source's
population-resolving motivation. It is not an impossibility statement for
one-fluid descriptions with memory, extra variables, gradients or nonlocality.
